% !TEX root = ../main.tex

\chapter{Analyse adaptiver
Lernverfahren}\label{analyse-adaptiver-lernverfahren}

\section{Forschungsstand und verwandte
Arbeiten}\label{forschungsstand-und-verwandte-arbeiten}

Die neuere Literatur beschreibt adaptive Lernsysteme nicht als einzelne,
klar umrissene Technologie, sondern als breites Feld personalisierter
digitaler Lernumgebungen, die Inhalte, Rückmeldungen, Aufgaben oder
Lernpfade auf Grundlage von Lernerdaten dynamisch anpassen
\citep{duPlooy2024, Wang2025, Martin2020}. Überblicksarbeiten
verdeutlichen dabei, dass adaptive Lernsysteme in sehr unterschiedlichen
Bildungskontexten - etwa in Hochschulumgebungen, intelligenten
Tutorensystemen oder digitalen Lernplattformen - eingesetzt werden und
zugleich methodisch stark ausdifferenzViert sind. Das Spektrum reicht von
klassischen psychometrischen und probabilistischen Modellen über
regelbasierte und wissensbasierte Verfahren bis hin zu
Recommender-Systemen sowie neueren KI- und Machine-Learning-Ansätzen
\citep{Wang2025, Hariyanto2025, daSilva2023, Liu2024}.\\
Adaptive Lerntechnologien lassen sich daher nicht auf einzelne
technische Ansätze reduzieren, sondern umfassen unterschiedliche
Verfahren, Modelle und Formen der Personalisierung.

Allgemeine Überblicksarbeiten betonen häufig die pädagogischen,
organisatorischen und implementierungsbezogenen Potenziale adaptiver
Lernsysteme. Für den Hochschulbereich zeigt ein aktueller Scoping
Review, dass adaptive Lernangebote häufig mit Verbesserungen von
Lernleistung und Engagement verbunden sind, zugleich jedoch durch
technische, zeitliche und organisatorische Ressourcen begrenzt werden.
Eine weitere Review zu KI-gestützten Adaptive-Learning-Plattformen hebt
hervor, dass insbesondere Plattformdesign, pädagogische Fundierung,
praktische Implementierung sowie geeignete Evaluationsmetriken zentrale
Themen des Forschungsfeldes darstellen. Solche Arbeiten machen die
Breite und Vielfalt adaptiver Lernsysteme deutlich, betrachten Verfahren
jedoch häufig primär auf Ebene konkreter Plattformen oder
Anwendungsszenarien. Die zugrunde liegenden algorithmischen Verfahren
werden dabei meist nicht systematisch als eigenständige
Verfahrensfamilien analysiert \citep{duPlooy2024, Tan2025}.

Daneben existieren zahlreiche spezialisierte Reviews, die jeweils nur
einzelne Teilbereiche adaptiver Lernsysteme untersuchen. Wang et
al.~analysieren adaptive E-Learning-Systeme mit Fokus auf die
Modellierung von Lernendenmerkmalen und die Verfahren zu deren
Identifikation. Die Autoren zeigen dabei insbesondere die große
methodische Vielfalt bestehender Learner-Modeling-Ansätze sowie eine
zunehmende Erweiterung des eingesetzten Methodenspektrums von
klassischen regelbasierten und probabilistischen Verfahren hin zu
breiteren Machine-Learning-Ansätzen \citep{Wang2025}. Desmarais und
Baker heben insbesondere probabilistische Modelle und Modelle latenter
Wissenszustände als zentrale Ansätze des Learner- und Skill-Modelings in
intelligenten Lernumgebungen hervor und zeigen, dass Verfahren zur
Wissensdiagnose eine Schlüsselrolle innerhalb adaptiver Systeme
einnehmen \citep{Desmarais2012}. Neuere Surveys zu Knowledge Tracing
beschreiben zudem einen Übergang von klassischen probabilistischen
Verfahren wie Bayesian Knowledge Tracing hin zu tiefen neuronalen
Sequenzmodellen \citep{Abdelrahman2023}. Im Bereich des Computerized
Adaptive Testing liegt der Fokus insbesondere auf Verfahren zur
adaptiven Auswahl informativer Testaufgaben sowie auf der
kontinuierlichen Fähigkeitsdiagnostik unter Verwendung psychometrischer
Modelle \citep{Han2018, Liu2024}. Aktuelle Reviews zu Reinforcement
Learning im Bildungsbereich zeigen, dass RL-Verfahren insbesondere zur
sequentiellen Optimierung adaptiver Lernentscheidungen eingesetzt
werden, zugleich jedoch weiterhin Herausforderungen hinsichtlich
realistischer Trainingsdaten und praktischer Implementierung bestehen
\citep{Riedmann2025}.\\
Der Forschungsstand ist damit reich an spezialisierten Teilanalysen,
enthält jedoch bislang nur wenige übergreifende Einordnungen, die
unterschiedliche Verfahrensfamilien systematisch und aus informatischer
Perspektive über Anwendungskontexte hinweg zusammenführen.

Darüber hinaus untersucht ein Teil der Literatur adaptive Lernsysteme
aus Perspektive von Implementierung, Akzeptanz und
Stakeholderbeteiligung. Alajlani, Crabb und Murray zeigen in ihrem
systematischen Review, dass Lehrende bislang nur in begrenztem Umfang in
Design, Implementierung und Evaluation adaptiver Lernsysteme eingebunden
sind \citep{Alajlani2024}. Diese Perspektive ist insbesondere für die
praktische Einführung adaptiver Lernsysteme relevant, trägt jedoch nur
eingeschränkt dazu bei, die zugrunde liegenden Verfahren und
Adaptionsmechanismen systematisch einzuordnen.

Insgesamt zeigt der Forschungsstand damit eine deutliche
Perspektivenpluralität: Pädagogische, didaktische, organisationale,
anwendungsbezogene und algorithmische Zugänge existieren weitgehend
nebeneinander, werden jedoch bislang nur selten in übergreifenden
Systematisierungen zusammengeführt.

\section{Methodisches Vorgehen des Scoping
Reviews}\label{methodisches-vorgehen-des-scoping-reviews}

Zur Analyse bestehender Klassifikationsansätze und algorithmischer
Verfahren adaptiver Lernsysteme wurde ein explorativ angelegtes Scoping
Review durchgeführt. Ziel war dabei nicht die vollständige Erfassung
sämtlicher Veröffentlichungen des Forschungsfeldes, sondern die
strukturierende Analyse zentraler Forschungsrichtungen, algorithmischer
Verfahren und bestehender Systematisierungsansätze.

Das Vorgehen orientierte sich an methodischen Empfehlungen für Scoping
Reviews nach Arksey und O'Malley sowie den weiterentwickelten Leitlinien
von Peters et al.~und PRISMA-ScR
\citep{Arksey2005, Peters2020, Tricco2018}. Übernommen wurden
insbesondere die explorative Ausrichtung des Vorgehens, die
systematische Identifikation relevanter Literatur sowie die thematische
Zusammenführung unterschiedlicher Forschungsstränge.\\
Auf eine vollständige PRISMA-konforme Durchführung mit formalisiertem
Screening-Prozess, detaillierter Dokumentation aller Suchanfragen und
vollständigem PRISMA-Flow-Diagramm wurde hingegen verzichtet.
Hintergrund hierfür ist, dass die Arbeit nicht auf eine vollständige
quantitative Erfassung sämtlicher Veröffentlichungen abzielt, sondern
auf die strukturierende Analyse zentraler algorithmischer Verfahren und
Klassifikationsperspektiven adaptiver Lernsysteme.

Die Literaturrecherche erfolgte explorativ unter Nutzung
wissenschaftlicher Suchmaschinen und Datenbanken. Ergänzend wurden
KI-gestützte Recherche- und Analysewerkzeuge verwendet, um relevante
Überblicksarbeiten, Surveys und thematische Zusammenhänge innerhalb des
Forschungsfeldes zu identifizieren. Berücksichtigt wurden insbesondere
Überblicksarbeiten, Surveys und systematische Reviews zu adaptiven
Lernsystemen, algorithmischen Verfahren adaptiver Lernprozesse, Learner
Modeling, Adaptivitätskonzepten, Empfehlungssystemen sowie bestehenden
Klassifikations- und Strukturierungsansätzen des Forschungsfeldes.\\
Die Auswahl der Literatur erfolgte anhand ihrer thematischen Relevanz
für die Forschungsfragen der Arbeit. Besonderes Augenmerk lag auf
Arbeiten, die bestehende Klassifikationen adaptiver Systeme beschreiben
oder algorithmische Verfahren und deren funktionale Rollen innerhalb
adaptiver Lernprozesse analysieren.\\
Für die weitere Analyse wurden insbesondere Überblicksarbeiten, Surveys
und systematische Reviews berücksichtigt, da diese eine vergleichende
Einordnung bestehender Verfahren und Forschungsperspektiven ermöglichen.
Die identifizierten Arbeiten wurden anschließend hinsichtlich ihrer
Klassifikationslogiken, ihrer funktionalen Perspektiven auf Adaptivität
sowie der beschriebenen algorithmischen Verfahren ausgewertet und
thematisch zusammengeführt.

\section{Bestehende Klassifikationen und
Perspektiven}\label{bestehende-klassifikationen-und-perspektiven}

Die Klassifikation adaptiver Lernsysteme gestaltet sich schwierig, da
die Literatur unterschiedliche Perspektiven auf Adaptivität und
Personalisierung verwendet. Während einige Arbeiten adaptive Systeme
primär über ihre Systemarchitektur beschreiben, orientieren sich andere
an pädagogischen Zielgrößen, psychologischen Merkmalen oder spezifischen
algorithmischen Teilproblemen. Eine einheitliche und allgemein
akzeptierte Taxonomie adaptiver Lernverfahren existiert bislang nicht.

Für die vorliegende Arbeit ist diese Heterogenität besonders relevant,
da unterschiedliche Klassifikationsansätze jeweils verschiedene
funktionale Aspekte adaptiver Systeme hervorheben. Im Folgenden werden
daher zentrale Perspektiven bestehender Systematisierungen analysiert
und hinsichtlich ihrer Eignung für eine informatisch orientierte
Klassifikation algorithmischer Verfahren eingeordnet.\\
Die Heterogenität der Literatur macht zugleich deutlich, dass eine
informatisch orientierte Klassifikation adaptiver Lernverfahren nicht
unmittelbar aus einer einzelnen bestehenden Taxonomie übernommen werden
kann. Vielmehr ergibt sich die Notwendigkeit, unterschiedliche
etablierte Perspektiven vergleichend zu betrachten und zu einer
begründeten Synthese zusammenzuführen.

\subsection{Adaptivitätsorientierte
Klassifikationen}\label{adaptivituxe4tsorientierte-klassifikationen}

Eine adaptivitätsorientierte Perspektive beschreibt adaptive Lernsysteme
danach, auf welche Arten von Lernendenmerkmalen Anpassungen erfolgen.
Plass und Pawar strukturieren Adaptivität dabei anhand affektiver,
kognitiver, motivationaler sowie soziokultureller Variablen
\citep{Plass2020}.

Zusätzlich unterscheiden Plass und Pawar zwischen Mikro- und
Makroadaptivität. Makroadaptivität beschreibt längerfristige oder
stärker vorausgeplante Anpassungen innerhalb des allgemeinen
Lernkontexts. Hierzu zählen beispielsweise die Gestaltung von
Lernpfaden, die Sequenzierung von Kursinhalten, vorbereitende
Lernaktivitäten, Unterstützungsangebote oder curriculare Entscheidungen
innerhalb eines Studienverlaufs. Anpassungen erfolgen dabei
typischerweise auf Grundlage allgemeinerer Informationen über Lernende
und können daher auch qualitative Merkmale und längerfristige Lernziele
berücksichtigen.

Mikroadaptivität bezieht sich dagegen auf kurzfristige und dynamische
Anpassungen innerhalb der unmittelbaren Lerninteraktion. Dazu gehören
beispielsweise adaptive Rückmeldungen, Hilfestellungen,
Aufgabenanpassungen oder die Auswahl nächster Lernschritte während der
Bearbeitung konkreter Lernaufgaben. Solche Anpassungen erfordern häufig
eine kontinuierliche oder echtzeitnahe Erfassung von Interaktionsdaten
und beruhen daher oft auf quantitativen Messungen des Lernverhaltens
\citep{Plass2020}.

Die Stärke dieser Taxonomie liegt insbesondere darin, dass Adaptivität
nicht primär technisch, sondern funktional über die zugrunde liegenden
Anpassungsziele beschrieben wird. Dadurch eignet sich der Ansatz gut, um
unterschiedliche Formen personalisierten Lernens pädagogisch und
psychologisch vergleichbar zu machen.

Für eine informatisch orientierte Analyse algorithmischer Verfahren
bleibt diese Perspektive jedoch nur begrenzt ausreichend. Die Taxonomie
fokussiert damit primär die funktionale Perspektive der Adaptivität,
weniger jedoch die algorithmischen Mechanismen, durch die adaptive
Entscheidungen tatsächlich realisiert werden.

\subsection{System- und lernendenmodellorientierte
Perspektiven}\label{system--und-lernendenmodellorientierte-perspektiven}

Im Unterschied zu adaptivitätsorientierten Ansätzen beschreiben andere
Arbeiten adaptive Lernsysteme stärker über ihre Systemarchitektur und
die Repräsentation lernendenbezogener Informationen. Typischerweise
werden dabei Komponenten wie Domänenmodell, Learner Model,
Adaptionsmodell sowie Benutzerschnittstelle unterschieden. Das
Domänenmodell repräsentiert dabei die fachlichen Inhalte und
Wissensstrukturen eines Lernsystems, während das Learner Model
Informationen über Wissen, Fähigkeiten, Präferenzen oder Interaktionen
von Lernenden enthält. Das Adaptionsmodell nutzt diese Informationen, um
Entscheidungen über Lerninhalte, Unterstützungsmaßnahmen oder Lernpfade
zu treffen. Die Benutzerschnittstelle bildet schließlich die Interaktion
zwischen Lernenden und System ab. Innerhalb dieser Architektur nimmt
insbesondere das Learner Model eine zentrale Rolle ein, da dort
diejenigen Informationen über Lernende repräsentiert werden, auf deren
Grundlage spätere adaptive Entscheidungen erfolgen \citep{Wang2025}.

Im Mittelpunkt dieser Perspektive stehen daher insbesondere Learner
Models sowie die Frage, wie Informationen über Lernende innerhalb
adaptiver Systeme strukturiert, verarbeitet und für
Adaptionsentscheidungen genutzt werden.

Martin et al.~verfolgen eine stärker forschungs- und
implementierungsorientierte Perspektive auf adaptive Lernsysteme. Statt
adaptive Verfahren primär algorithmisch zu klassifizieren,
systematisieren die Autoren die Forschungslandschaft anhand von
Publikationstrends, Anwendungskontexten, eingesetzten Strategien und
technologischen Umsetzungen adaptiver Systeme.\\
Die Analyse zeigt dabei unter anderem, dass Lernstile über längere Zeit
zu den am häufigsten betrachteten Lernendenmerkmalen gehörten.
Gleichzeitig wurden insbesondere adaptive Feedback-Mechanismen und
adaptive Navigation häufig untersucht. Darüber hinaus verdeutlicht die
Arbeit, dass adaptive Lernsysteme in sehr unterschiedlichen
Bildungskontexten eingesetzt werden und sich sowohl hinsichtlich ihrer
pädagogischen Zielsetzungen als auch ihrer technischen Umsetzung stark
unterscheiden \citep{Martin2020}.

Auch Wang et al.~verfolgen keine primär algorithmische Klassifikation
adaptiver Systeme, sondern betrachten adaptive Lernumgebungen aus
Perspektive des Learner Modelings. Im Mittelpunkt steht dabei die Frage,
welche Lernendenmerkmale innerhalb adaptiver Systeme erfasst und
modelliert werden und mit welchen Verfahren diese Merkmale identifiziert
oder aktualisiert werden. Die Autoren beschreiben adaptive Systeme dabei
als mehrstufigen Prozess: Zunächst werden relevante Merkmale von
Lernenden -- beispielsweise Wissen, Präferenzen oder Lernverhalten --
innerhalb eines Learner Models repräsentiert. Anschließend kommen
unterschiedliche Verfahren zum Einsatz, um Lernende anhand dieser
Merkmale zu klassifizieren oder Veränderungen des Lernzustands zu
erkennen. Auf Grundlage dieser Informationen erfolgen schließlich
adaptive Entscheidungen, etwa bei der Auswahl von Lerninhalten, Aufgaben
oder Unterstützungsmaßnahmen \citep{Wang2025}.

Eine noch stärkere Fokussierung auf das Learner Model selbst findet sich
bei Böck. Dort wird das Learner Model als explizite maschineninterne
Repräsentation individueller Merkmale, Zustände und Interaktionsdaten
eines Lernenden beschrieben. Im Mittelpunkt steht dabei insbesondere die
Frage, wie lernendenbezogene Informationen innerhalb adaptiver Systeme
formal strukturiert, repräsentiert und weiterverarbeitet werden.\\
Zugleich zeigt die Arbeit, dass viele Publikationen zwar mögliche
Attribute von Learner Models benennen, deren konkrete formale
Modellierung sowie die zugrunde liegenden Datenquellen jedoch häufig
unklar bleiben. Darüber hinaus verweist Böck darauf, dass bislang kein
allgemeiner und flexibel erweiterbarer Standard für Learner Models
identifiziert werden konnte.~ Die fehlende Standardisierung von Learner
Models hängt zugleich mit der großen Vielfalt unterschiedlicher
Modellierungsansätze zusammen. Böck verweist in diesem Zusammenhang auf
Verfahren, die von regel- und constraintsbasierten Ansätzen über
probabilistische Verfahren und Ontologien bis hin zu Clustering-,
Klassifikations- und Deep-Learning-Methoden reichen. Dies verdeutlicht,
dass adaptive Lernsysteme nicht nur unterschiedliche Lernendenmerkmale
modellieren, sondern hierfür auch sehr unterschiedliche technische
Modellierungsstrategien verwenden \citep{Boeck2025}.

Die betrachteten Arbeiten verdeutlichen insgesamt unterschiedliche
Perspektiven auf adaptive Lernsysteme. Während Martin et al.~adaptive
Lernforschung primär über Kontexte, Strategien und Technologien
strukturieren, stellen Wang et al.~insbesondere den Prozess des Learner
Modelings in den Mittelpunkt. Böck fokussiert dagegen stärker die
technische Struktur und Modellierung lernendenbezogener Informationen.
Gemeinsam ist diesen Ansätzen, dass sie wichtige Einblicke in Aufbau,
Informationsverarbeitung und Adaptionslogik adaptiver Systeme liefern,
jedoch keine durchgängige Einordnung algorithmischer Verfahrensfamilien
vornehmen. Adaptive Lernsysteme werden damit häufig über
Lernendenmerkmale, Systemarchitekturen oder Forschungs- und
Anwendungskontexte beschrieben, weniger jedoch über die zugrunde
liegenden algorithmischen Mechanismen selbst.

\subsection{Algorithmisch orientierte
Perspektiven}\label{algorithmisch-orientierte-perspektiven}

Im Gegensatz zu stärker pädagogisch, system- oder
lernendenmodellorientierten Ansätzen existieren auch Arbeiten, die
adaptive Lernsysteme primär über algorithmische Verfahren und Methoden
der Wissensmodellierung strukturieren. Im Mittelpunkt steht dabei
weniger die Frage, welche Lernendenmerkmale modelliert werden, sondern
vielmehr, mit welchen technischen Verfahren adaptive Entscheidungen
getroffen, Wissenszustände geschätzt oder Lernprozesse analysiert
werden.

Zwar existieren zahlreiche spezialisierte Arbeiten zu einzelnen
algorithmischen Verfahrensfamilien adaptiver Lernsysteme, darunter
probabilistische Verfahren des Learner Modelings, psychometrische
Modelle im Computerized Adaptive Testing, Verfahren des Knowledge
Tracing, Recommender-Systeme oder Reinforcement-Learning-basierte
Ansätze. Diese Arbeiten untersuchen adaptive Verfahren jedoch meist
innerhalb spezifischer Anwendungskontexte oder Forschungsfelder und
verwenden unterschiedliche Begrifflichkeiten, Modellannahmen und
Klassifikationslogiken. Eine übergreifende und konsistente
algorithmische Taxonomie adaptiver Lernverfahren findet sich bislang nur
eingeschränkt.

Im Bereich des Learner Modelings spielen insbesondere probabilistische
Verfahren zur Wissens- und Fähigkeitsdiagnostik eine zentrale Rolle.
Desmarais und Baker beschreiben verschiedene Ansätze zur Modellierung
von Wissenszuständen in intelligenten Lernumgebungen und heben hervor,
dass probabilistische Modelle besonders geeignet sind, Unsicherheit
innerhalb adaptiver Systeme abzubilden \citep{Desmarais2012}. Adaptive
Systeme besitzen dabei keinen direkten Zugriff auf den tatsächlichen
Wissensstand von Lernenden, sondern müssen diesen aus beobachtbaren
Interaktionen wie Antworten, Fehlermustern oder Bearbeitungszeiten
inferieren.

Pelánek beschreibt in seiner Überblicksarbeit zum Learner Modeling
insbesondere Bayesian Knowledge Tracing sowie logistische Modelle als
zentrale Ansätze zur Modellierung von Wissenszuständen über mehrere
Lerninteraktionen hinweg \citep{Pelanek2017}. Gleichzeitig betont die
Arbeit, dass unterschiedliche Modellierungsverfahren jeweils von den
verfügbaren Daten, dem Anwendungskontext und dem Ziel der Adaptation
abhängen. Lernendenmodellierung stellt damit keinen einzelnen
Algorithmus dar, sondern einen breiten Bereich unterschiedlicher
Zustands-, Fähigkeits- und Wahrscheinlichkeitsmodelle.

Neuere Arbeiten zum Knowledge Tracing differenzieren adaptive
Wissensmodellierung zusätzlich stärker nach technischen Modellfamilien.
Abdelrahman et al.~unterscheiden dabei insbesondere Bayesian Models,
Logistic Models sowie Deep-Learning-basierte Ansätze
\citep{Abdelrahman2023}. Neben klassischen Verfahren wie Bayesian
Knowledge Tracing treten damit zunehmend neuronale Sequenzmodelle zur
Modellierung latenter Wissenszustände in den Vordergrund. Die Literatur
verdeutlicht zugleich, dass ähnliche funktionale Ziele -- etwa die
Schätzung von Wissensständen -- mit sehr unterschiedlichen
Modellierungsansätzen realisiert werden können.

Neben probabilistischen Verfahren existieren weitere spezialisierte
algorithmische Forschungsfelder adaptiver Lernsysteme. Im Bereich des
Computerized Adaptive Testing stehen insbesondere psychometrische
Verfahren zur adaptiven Auswahl informativer Testaufgaben und zur
kontinuierlichen Fähigkeitsdiagnostik im Mittelpunkt
\citep{Han2018, Liu2024}. Educational Recommender Systems konzentrieren
sich dagegen auf die personalisierte Auswahl von Lernressourcen,
Aufgaben oder Lernpfaden auf Grundlage von Präferenzen,
Interaktionsdaten oder Ähnlichkeiten zwischen Lernenden
\citep{daSilva2023, Nabizadeh2020}. Reinforcement-Learning-basierte
Ansätze fokussieren schließlich die sequenzielle Optimierung adaptiver
Entscheidungen und die dynamische Steuerung von Lernprozessen
\citep{Riedmann2025}.

Die algorithmisch orientierte Perspektive macht deutlich, dass adaptive
Lernsysteme eine große Vielfalt unterschiedlicher Modellierungs- und
Entscheidungsverfahren verwenden. Gleichzeitig zeigt sich jedoch, dass
die Literatur algorithmische Verfahren häufig innerhalb spezialisierter
Teilbereiche diskutiert und nur selten übergreifend zusammenführt. Die
bestehende Forschungslandschaft ist damit reich an spezialisierten
Verfahrensanalysen, weist jedoch weiterhin eine starke Fragmentierung
hinsichtlich Begrifflichkeiten, Modellierungslogiken und
Klassifikationsansätzen auf.

\subsection{KI- und Machine-Learning-orientierte
Perspektiven}\label{ki--und-machine-learning-orientierte-perspektiven}

Neben probabilistischen, psychometrischen oder regelbasierten Verfahren
existieren auch Arbeiten, die adaptive Lernsysteme primär über Methoden
des maschinellen Lernens strukturieren. Im Mittelpunkt steht dabei
insbesondere die automatische Analyse von Lernerdaten sowie die
datengetriebene Modellierung adaptiver Entscheidungen.

Hariyanto et al.~klassifizieren adaptive Bildungssysteme beispielsweise
anhand zentraler Kategorien des maschinellen Lernens und unterscheiden
zwischen überwachten, unüberwachten und
Reinforcement-Learning-Verfahren.

Überwachte Lernverfahren werden dabei vor allem zur Klassifikation von
Lernendenprofilen und zur Vorhersage lernbezogener Eigenschaften
eingesetzt. Die Autoren nennen unter anderem Entscheidungsbäume, Support
Vector Machines sowie k-Nearest-Neighbor-Verfahren als typische Ansätze
zur Analyse von Leistungsdaten und Lernverhalten. Ziel solcher Verfahren
ist es häufig, Lernende bestimmten Kategorien oder
Unterstützungsmaßnahmen zuzuordnen.

Unüberwachte Lernverfahren dienen dagegen insbesondere der
Identifikation von Mustern und Gruppenstrukturen innerhalb von
Lernerdaten. Hierzu werden beispielsweise Clustering-Verfahren genutzt,
um Ähnlichkeiten in Interaktions-, Zeit- oder Leistungsdaten zu erkennen
und Lernende anhand gemeinsamer Merkmale zu gruppieren.

Reinforcement-Learning-Verfahren werden vor allem zur dynamischen
Steuerung adaptiver Lernprozesse eingesetzt. Reinforcement Learning
ermöglicht dabei insbesondere die sequenzielle Optimierung von
Lernpfaden, Inhaltsauswahl oder Unterstützungsmaßnahmen auf Grundlage
fortlaufender Rückmeldungen aus der Interaktion mit dem Lernsystem
\citep{Hariyanto2025}.

Die KI- und ML-orientierte Perspektive verdeutlicht insbesondere die
zunehmende Bedeutung datengetriebener Verfahren innerhalb adaptiver
Lernsysteme. Gleichzeitig konzentrieren sich solche Arbeiten häufig auf
bestimmte Paradigmen des maschinellen Lernens, während probabilistische,
psychometrische oder symbolische Verfahren teilweise weniger
berücksichtigt werden.

\subsection{Zwischenbetrachtung und Grenzen bestehender
Ansätze}\label{zwischenbetrachtung-und-grenzen-bestehender-ansuxe4tze}

Die analysierten Arbeiten verdeutlichen insgesamt, dass adaptive
Lernsysteme in der Forschung aus sehr unterschiedlichen Perspektiven
beschrieben und strukturiert werden. Während einige Ansätze Adaptivität
primär über Lernendenmerkmale und pädagogische Zielgrößen definieren,
fokussieren andere stärker Systemarchitekturen, Learner Models,
Forschungs- und Anwendungskontexte oder konkrete algorithmische
Verfahren.

Darüber hinaus zeigt sich, dass adaptive Lernsysteme häufig innerhalb
spezialisierter Forschungsfelder wie Intelligent Tutoring Systems,
Knowledge Tracing, Educational Recommender Systems, Computerized
Adaptive Testing oder Reinforcement-Learning-basierten Lernumgebungen
untersucht werden
\citep{Desmarais2012, Abdelrahman2023, daSilva2023, Liu2024, Riedmann2025}.
Dadurch entstehen eigenständige Teilbereiche mit jeweils eigenen
Begrifflichkeiten, Modellierungsansätzen und Bewertungskriterien.

Gleichzeitig wird jedoch deutlich, dass ähnliche algorithmische
Grundprinzipien in unterschiedlichen Anwendungskontexten wiederkehren.
Probabilistische Zustandsmodellierung, Klassifikation, Mustererkennung
oder sequenzielle Entscheidungsoptimierung treten beispielsweise in
mehreren Teilbereichen adaptiver Lernsysteme auf, werden jedoch häufig
getrennt voneinander betrachtet und nur selten gemeinsam eingeordnet.

Die bestehende Literatur liefert damit zahlreiche spezialisierte
Teilsystematiken und wichtige Perspektiven auf adaptive Lernsysteme,
bietet bislang jedoch nur begrenzte Ansätze für eine integrierte
informatisch orientierte Klassifikation algorithmischer Verfahren.
Insbesondere fehlt häufig eine gemeinsame Strukturierung, die
funktionale Rolle, Datenbasis, Adaptionsmechanismus und technische
Modellierungsansätze systematisch zusammenführt.

\section{Zentrale algorithmische Verfahren im
Überblick}\label{zentrale-algorithmische-verfahren-im-uxfcberblick}

Für eine informatisch orientierte Betrachtung adaptiver Lernsysteme ist
zunächst weniger entscheidend, welcher konkrete Plattformtyp vorliegt,
sondern vielmehr, welche algorithmischen Verfahrensfamilien innerhalb
adaptiver Systeme wiederkehren und welche funktionalen Rollen diese
übernehmen.

\subsection{Learner Modeling und Knowledge
Tracing}\label{learner-modeling-und-knowledge-tracing}

Eine zentrale Rolle spielt dabei der Bereich des Learner Modeling
beziehungsweise der Wissens- und Fähigkeitsdiagnostik.

Learner Modeling beschäftigt sich grundsätzlich mit der Frage, wie
Wissen, Fähigkeiten, Präferenzen oder weitere lernendenbezogene Merkmale
innerhalb eines adaptiven Systems repräsentiert und fortlaufend
aktualisiert werden können \citep{Wang2025, Boeck2025}. Adaptive Systeme
besitzen dabei keinen direkten Zugriff auf den tatsächlichen
Wissensstand eines Lernenden. Beobachtbar sind lediglich Interaktionen
mit dem System, etwa Antworten auf Aufgaben, Fehlversuche,
Bearbeitungszeiten oder Navigationsverhalten. Der eigentliche
Wissenszustand stellt somit eine latente Größe dar, die aus
beobachtbaren Daten erschlossen werden muss \citep{Desmarais2012}.

Vor diesem Hintergrund spielen probabilistische Verfahren eine zentrale
Rolle im Learner Modeling. Desmarais und Baker betonen, dass
probabilistische Modelle für die Einschätzung von Fähigkeiten und
Wissenszuständen in intelligenten Lernumgebungen von grundlegender
Bedeutung sind. Solche Verfahren gehen davon aus, dass Wissen nicht als
direkt beobachtbarer Zustand vorliegt, sondern nur mit einer bestimmten
Wahrscheinlichkeit angenommen werden kann. Anstatt einen Lernenden
eindeutig als \enquote{beherrscht} oder \enquote{nicht beherrscht} zu
klassifizieren, modellieren probabilistische Ansätze daher die
Wahrscheinlichkeit, mit der bestimmte Kenntnisse oder Fähigkeiten
vorhanden sind. Neue Interaktionen mit dem Lernsystem können
anschließend genutzt werden, um diese Wahrscheinlichkeiten fortlaufend
anzupassen und die Einschätzung des Wissenszustands zu aktualisieren
\citep{Desmarais2012}.

In der späteren Überblicksliteratur zum Learner Modeling hebt Pelánek
insbesondere Bayesian Knowledge Tracing sowie logistische Modelle als
wichtige Ansätze zur Modellierung von Wissenszuständen hervor. Bayesian
Knowledge Tracing modelliert dabei den Wissenszustand eines Lernenden
als Wahrscheinlichkeit, die nach jeder bearbeiteten Aufgabe aktualisiert
wird. Logistische Modelle verfolgen dagegen einen stärker statistischen
Ansatz und schätzen die Wahrscheinlichkeit korrekter Antworten auf
Grundlage bestimmter Merkmale von Lernenden oder Aufgaben. Beide Ansätze
verfolgen damit das Ziel, Lernprozesse quantitativ zu beschreiben und
zukünftige Lernleistungen vorherzusagen, beruhen jedoch auf
unterschiedlichen Modellannahmen über Wissenserwerb und
Kompetenzentwicklung \citep{Pelanek2017}.\\
Gleichzeitig betont Pelánek, dass die Wahl eines geeigneten Modells
stets vom jeweiligen Anwendungskontext, den verfügbaren Daten sowie dem
Ziel der Adaptation abhängt \citep{Pelanek2017}. Lernendenmodellierung
stellt damit keinen einzelnen Algorithmus dar, sondern einen breiten
Bereich unterschiedlicher Zustands-, Fähigkeits- und
Wahrscheinlichkeitsmodelle, die jeweils verschiedene Perspektiven auf
Lernen, Wissen und Kompetenzentwicklung abbilden.

Besonders stark ausdifferenziert ist dieser Bereich im Forschungsfeld
des Knowledge Tracing. Knowledge Tracing beschäftigt sich mit der Frage,
wie sich der Wissenszustand eines Lernenden über mehrere
Lerninteraktionen hinweg modellieren und fortlaufend aktualisieren
lässt. Ziel ist es, auf Grundlage bisheriger Bearbeitungen möglichst
präzise abzuschätzen, welche Wissenskomponenten bereits beherrscht
werden und bei welchen weiterhin Lernbedarf besteht
\citep{Abdelrahman2023}.

Der Survey von Shen et al.~unterscheidet drei grundlegende technische
Modellierungsrichtungen innerhalb des Knowledge Tracing: Bayesian
Models, Logistic Models und Deep Learning Models \citep{Shen2024}. Diese
Einteilung ist für eine informatische Systematisierung besonders
relevant, da die Verfahren nicht primär nach ihrem Anwendungskontext,
sondern nach ihrer technischen Modellierungslogik gruppiert werden.

Bayesian Models modellieren Wissen als verborgenen beziehungsweise
latenten Zustand, dessen Wahrscheinlichkeit anhand beobachtbarer
Interaktionen geschätzt wird. Zu den bekanntesten Vertretern zählen
Bayesian Knowledge Tracing (BKT) sowie Erweiterungen wie Dynamic
Bayesian Knowledge Tracing (DBKT). Der Wissenszustand eines Lernenden
wird dabei als Wahrscheinlichkeit dargestellt, die nach jeder
bearbeiteten Aufgabe auf Grundlage neuer Beobachtungen aktualisiert wird
\citep{Pelanek2017, Shen2024}.

Logistic Models verfolgen dagegen einen stärker statistischen Ansatz. Zu
den wichtigsten Vertretern zählen Learning Factors Analysis (LFA) und
Performance Factors Analysis (PFA). Im Mittelpunkt steht hierbei die
Schätzung der Wahrscheinlichkeit korrekter Antworten auf Grundlage
beobachtbarer Merkmale und Interaktionsdaten, etwa früherer Erfolge oder
Fehlversuche. Im Unterschied zu klassischen Bayesian-Modellen wird dabei
nicht unmittelbar ein verborgener Wissenszustand modelliert, sondern das
Antwortverhalten selbst statistisch beschrieben \citep{Shen2024}.

Deep-Learning-basierte Verfahren verzichten weitgehend auf manuell
definierte Wissensmodelle und lernen Zusammenhänge direkt aus großen
Mengen von Interaktionsdaten. Zu den bekanntesten Ansätzen zählt Deep
Knowledge Tracing (DKT), das rekurrente neuronale Netze zur Modellierung
von Lernverläufen verwendet. Neuere Varianten erweitern diesen Ansatz
beispielsweise um Speichermechanismen, Aufmerksamkeitsstrukturen oder
Graphrepräsentationen. Dadurch können komplexe Muster im Lernverhalten
erfasst werden, allerdings geht dies häufig zulasten der
Interpretierbarkeit der Modelle \citep{Shen2024}.

Die von Shen et al.~vorgeschlagene Einteilung verdeutlicht, dass
dieselbe funktionale Aufgabe -- die Schätzung von Wissenszuständen --
mit sehr unterschiedlichen technischen Modellierungsansätzen realisiert
werden kann. Für eine informatische Klassifikation adaptiver Verfahren
sind daher nicht nur die funktionalen Ziele, sondern auch die zugrunde
liegenden Modellierungslogiken von zentraler Bedeutung.

Gleichzeitig zeigt die Knowledge-Tracing-Literatur, dass sich die
verschiedenen Verfahren nicht nur hinsichtlich ihrer
Modellierungsansätze, sondern auch hinsichtlich der verwendeten
Datenbasis unterscheiden. Die Qualität der Wissensschätzung hängt
wesentlich davon ab, welche Informationen einem Modell zur Verfügung
stehen und wie diese in die Schätzung des Lernzustands einbezogen werden
\citep{Shen2024}. Frühe Knowledge-Tracing-Modelle arbeiten häufig
ausschließlich mit Aufgabenfolgen und den zugehörigen Antwortdaten. Das
Modell beobachtet dabei typischerweise, ob eine Aufgabe korrekt oder
inkorrekt beantwortet wurde, und nutzt diese Informationen zur
Aktualisierung des geschätzten Wissenszustands. Solche Ansätze sind
vergleichsweise einfach zu implementieren und benötigen nur wenige
Eingabedaten, können jedoch viele Einflussfaktoren des tatsächlichen
Lernprozesses nicht berücksichtigen \citep{Shen2024}. Neuere Verfahren
erweitern diese Datenbasis deutlich. Neben den eigentlichen Antwortdaten
werden zunehmend zusätzliche Informationen in die Modellierung
einbezogen. Hierzu zählen beispielsweise individuelle Eigenschaften
eines Lernenden vor Beginn des Lernprozesses, Merkmale des bisherigen
Lernverlaufs, Indikatoren für Motivation oder Engagement sowie
Informationen über Vergessensprozesse im Zeitverlauf. Darüber hinaus
werden sogenannte Side Information genutzt, etwa Antwortzeiten,
Bearbeitungsdauer, Navigationsmuster oder Tutorinterventionen während
des Lernprozesses \citep{Shen2024}. Diese Entwicklung verdeutlicht, dass
dieselbe funktionale Aufgabe -- die Schätzung von Wissenszuständen --
mit sehr unterschiedlichen Datentiefegraden realisiert werden kann.
Während einfache Modelle lediglich auf einer Folge richtiger und
falscher Antworten basieren, versuchen komplexere Verfahren ein
umfassenderes Bild des Lernprozesses zu erfassen \citep{Shen2024}.

Für die Frage der Interpretierbarkeit ist diese Unterscheidung von
zentraler Bedeutung. Der Survey von Shen et al.~weist ausdrücklich
darauf hin, dass moderne Deep-Learning-Verfahren im Knowledge Tracing
häufig eine höhere Vorhersagegenauigkeit erreichen als traditionelle
Ansätze, ihre Entscheidungsprozesse jedoch deutlich schwerer
nachvollziehbar sind \citep{Shen2024}. Während klassische Modelle meist
auf klar definierten Parametern und Modellannahmen beruhen, entstehen
die Vorhersagen tiefer neuronaler Netze aus einer großen Zahl interner
Berechnungen, deren Einfluss auf das Ergebnis oft nur schwer zu
interpretieren ist.

Gerade im Bildungskontext besitzt diese Problematik besondere Relevanz.
Adaptive Lernsysteme treffen Entscheidungen über Lerninhalte,
Empfehlungen, Unterstützungsmaßnahmen oder die Einschätzung von
Wissensständen. Lehrende und Lernende haben dabei häufig ein
berechtigtes Interesse daran, nachvollziehen zu können, warum ein System
eine bestimmte Einschätzung vorgenommen oder eine konkrete Empfehlung
ausgesprochen hat. Fehlende Transparenz kann die Akzeptanz adaptiver
Systeme reduzieren und erschwert zugleich die pädagogische Bewertung
ihrer Entscheidungen.

Vor diesem Hintergrund hat sich innerhalb der
Knowledge-Tracing-Forschung ein eigener Forschungsstrang zu
interpretierbaren Modellen und Explainable Artificial Intelligence (XAI)
entwickelt. Ziel dieser Arbeiten ist es, die Entscheidungsprozesse
komplexer Modelle besser verständlich zu machen und zusätzliche
Erklärungen für Vorhersagen bereitzustellen. Dabei werden beispielsweise
Einflussfaktoren einzelner Aufgaben, Wissenskomponenten oder
Interaktionen sichtbar gemacht, um die Ursachen einer Modellentscheidung
nachvollziehbarer darzustellen \citep{Shen2024}.

Für eine informatische Klassifikation adaptiver Verfahren ergibt sich
daraus eine wichtige zusätzliche Betrachtungsdimension. Modelle
unterscheiden sich nicht nur hinsichtlich ihrer Vorhersageleistung oder
ihrer technischen Modellierungslogik, sondern auch hinsichtlich ihrer
Transparenz und Erklärbarkeit. Probabilistische und logistische
Verfahren sind dabei häufig leichter interpretierbar als
Deep-Learning-Ansätze, während tiefe neuronale Modelle oftmals eine
höhere Modellkomplexität und geringere Nachvollziehbarkeit aufweisen
\citep{Shen2024}. Die Interpretierbarkeit eines Verfahrens stellt daher
eine eigenständige Eigenschaft dar, die bei einer systematischen
Einordnung adaptiver Lernverfahren berücksichtigt werden sollte.

\subsection{Recommender-Systeme und
Lernpfadverfahren}\label{recommender-systeme-und-lernpfadverfahren}

Eine zweite große Verfahrensfamilie bilden Recommender-Systeme und
Lernpfadverfahren. Während Verfahren des Learner Modelings primär auf
die Diagnose und Modellierung von Wissenszuständen abzielen, steht bei
Recommender-Systemen die Auswahl und Priorisierung geeigneter
Lernressourcen im Vordergrund. Ziel ist es, Lernenden diejenigen
Inhalte, Aufgaben oder Lernaktivitäten bereitzustellen, die für ihre
individuellen Bedürfnisse, Interessen oder Lernziele besonders relevant
erscheinen \citep{daSilva2023}.

Im Bildungsbereich werden Recommender-Systeme daher als zentrales
Instrument zur Personalisierung von Lernumgebungen betrachtet.
Empfehlungen können sich beispielsweise auf Lernmaterialien, Aufgaben,
Kurse, Lernpfade oder unterstützende Ressourcen beziehen. Grundlage
hierfür bilden typischerweise Informationen über Lernende, deren
bisheriges Lernverhalten, Leistungsdaten, Präferenzen oder Ähnlichkeiten
zu anderen Nutzenden \citep{daSilva2023}.

Der systematische Review von da Silva et al.~zeigt, dass insbesondere
hybride Recommender-Systeme in der aktuellen Forschung und Praxis eine
dominante Rolle einnehmen. Hybride Verfahren kombinieren dabei
unterschiedliche Empfehlungsstrategien, etwa kollaboratives Filtern,
inhaltsbasierte Verfahren oder wissensbasierte Ansätze. Ziel dieser
Kombination ist es, die jeweiligen Stärken einzelner Methoden zu nutzen
und deren Schwächen auszugleichen. So können beispielsweise Empfehlungen
nicht nur auf Ähnlichkeiten zwischen Lernenden beruhen, sondern
zusätzlich fachliche Eigenschaften von Lernressourcen oder pädagogische
Anforderungen berücksichtigen \citep{daSilva2023}.

Die einzelnen Empfehlungsstrategien werden wiederum durch
unterschiedliche algorithmische Verfahren realisiert. Kollaborative
Recommender nutzen beispielsweise Ähnlichkeitsmaße oder
Nachbarschaftsverfahren wie k-Nearest Neighbors, um Lernende mit
vergleichbaren Interessen oder Verhaltensmustern zu identifizieren.
Inhaltsbasierte Ansätze arbeiten häufig mit Merkmalsvektoren und
Ähnlichkeitsmaßen wie der Kosinus-Ähnlichkeit, um passende
Lernressourcen zu bestimmen. Darüber hinaus kommen in Educational
Recommender Systems Verfahren wie Matrix Factorization, Fuzzy Logic,
genetische Algorithmen sowie zunehmend neuronale Netze und
Deep-Learning-Ansätze zum Einsatz \citep{daSilva2023}.

Gleichzeitig weist die Review darauf hin, dass die Evaluation von
Educational Recommender Systems bislang überwiegend auf technischen
Qualitätsmaßen wie Genauigkeit, Präzision oder Trefferquote basiert,
während die tatsächliche pädagogische Wirksamkeit deutlich seltener
untersucht wird \citep{daSilva2023}. Für die Analyse adaptiver
Lernverfahren ist dieser Befund besonders relevant, da
Recommender-Systeme nicht primär der Wissensdiagnostik dienen, sondern
der Auswahl und Priorisierung von Lernressourcen. Der Erfolg eines
Verfahrens lässt sich daher nicht allein an der Qualität der Empfehlung
messen, sondern sollte idealerweise auch anhand von Lernerfolg,
Motivation oder Lernerfahrung bewertet werden.

Innerhalb dieses Bereichs sind insbesondere wissens- und
ontologiebasierte Recommender-Systeme von besonderem Interesse. Im
Unterschied zu vielen datengetriebenen Verfahren beruhen sie nicht
primär auf statistischen Mustern oder Ähnlichkeiten zwischen Nutzenden,
sondern auf explizit modelliertem Wissen über Lernende, Lernressourcen
und fachliche Zusammenhänge. Die Personalisierung erfolgt dabei auf
Grundlage formalisierter Wissensstrukturen, welche Beziehungen zwischen
Konzepten, Lernzielen, Kompetenzen oder Inhaltsbereichen beschreiben
\citep{George2019}.

Ontologien dienen in diesem Zusammenhang als formale Repräsentationen
von Wissen. Sie definieren Konzepte eines bestimmten Anwendungsbereichs
sowie die Beziehungen zwischen diesen Konzepten. Dadurch können
Lerninhalte, Wissensgebiete, Kompetenzen und Lernendenmerkmale
strukturiert beschrieben und maschinell verarbeitet werden. Adaptive
Systeme können diese Wissensstrukturen nutzen, um Zusammenhänge zwischen
Lernenden und Lernressourcen abzuleiten und darauf aufbauend
personalisierte Empfehlungen zu erzeugen \citep{George2019}.

George und Lal heben hervor, dass ontologiebasierte Ansätze insbesondere
Reusability, Reasoning und Inferenz unterstützen \citep{George2019}. Da
Wissen explizit modelliert wird, können Systeme nicht nur vorhandene
Informationen speichern, sondern zusätzliche Beziehungen und
Schlussfolgerungen aus den modellierten Strukturen ableiten. Dadurch
eignen sich Ontologien besonders für adaptive Lernumgebungen, in denen
fachliche Abhängigkeiten, Kompetenzmodelle oder curriculare
Zusammenhänge berücksichtigt werden sollen.

Die Review von George und Lal zeigt zugleich, dass ontologiebasierte
Recommender-Systeme unterschiedliche technische Verfahren zur
Generierung von Empfehlungen einsetzen \citep{George2019}. Häufig kommen
dabei regelbasierte Schlussfolgerungsverfahren (Rule-Based Reasoning),
semantische Ähnlichkeitsmaße (Semantic Similarity), Case-Based Reasoning
sowie hybride Ansätze zum Einsatz, welche Ontologien mit kollaborativen
Filterverfahren, Fuzzy-Logic-Methoden oder weiteren Empfehlungstechniken
kombinieren. Die eigentliche Empfehlung entsteht dabei nicht allein
durch statistische Muster in den Nutzungsdaten, sondern durch
Schlussfolgerungen auf Grundlage der modellierten Wissensstrukturen.

Besonders semantische Ähnlichkeitsverfahren spielen eine wichtige Rolle.
Dabei werden Beziehungen zwischen Konzepten innerhalb einer Ontologie
genutzt, um die fachliche Nähe zwischen Lernressourcen, Kompetenzen oder
Lernzielen zu bestimmen. Regelbasierte Verfahren ermöglichen darüber
hinaus die Formulierung expliziter pädagogischer oder fachlicher Regeln,
auf deren Grundlage Empfehlungen abgeleitet werden können. Adaptive
Entscheidungen beruhen damit stärker auf symbolischen Inferenzprozessen
als auf rein datengetriebenem Lernen \citep{George2019}.

Gleichzeitig verweisen die Autoren auf mehrere Herausforderungen. Die
Erstellung und Pflege ontologischer Wissensmodelle ist häufig mit hohem
Modellierungsaufwand verbunden und erfordert umfangreiche
Domänenkenntnisse. Darüber hinaus wird die Evaluation solcher Systeme
durch das Fehlen standardisierter Datensätze und Vergleichsgrundlagen
erschwert \citep{George2019}.

Für eine informatische Klassifikation adaptiver Verfahren sind
ontologiebasierte Recommender insbesondere deshalb relevant, weil sie
eine deutlich andere Modellierungslogik verfolgen als rein
datengetriebene Ansätze. Während viele Machine-Learning-Verfahren
Zusammenhänge implizit aus Interaktionsdaten lernen, beruhen
ontologiebasierte Systeme auf explizit repräsentiertem Wissen und
symbolischen Schlussfolgerungsmechanismen. Adaptive Lernsysteme
unterscheiden sich daher nicht nur hinsichtlich ihrer Algorithmen,
sondern auch hinsichtlich der Art, wie Wissen innerhalb des Systems
repräsentiert und verarbeitet wird.

Verwandt, aber nicht deckungsgleich mit Recommender-Systemen sind
Lernpfad- und Sequenzierungsverfahren. Während klassische
Recommender-Systeme primär die Auswahl einzelner Lernressourcen oder
Inhalte unterstützen, beschäftigen sich Lernpfadverfahren mit der
Planung, Generierung und Anpassung ganzer Sequenzen von Lernaktivitäten.
Im Mittelpunkt steht dabei die Frage, in welcher Reihenfolge Lerninhalte
bearbeitet werden sollten, um individuelle Lernziele möglichst effizient
zu erreichen \citep{Nabizadeh2020}.

Der Survey von Nabizadeh et al.~unterscheidet zur Personalisierung von
Lernpfaden drei zentrale Ansätze: Course Generation, Sequential Pattern
Recognition und Course Sequence. Bei der Course Generation werden
individuelle Lernpfade auf Grundlage von Lernzielen, Vorwissen oder
Kompetenzanforderungen erzeugt. Sequential Pattern Recognition
konzentriert sich dagegen auf die Identifikation erfolgreicher
Lernsequenzen innerhalb vorhandener Lerndaten, um daraus geeignete
Reihenfolgen von Lernaktivitäten abzuleiten. Course-Sequence-Verfahren
fokussieren schließlich die gezielte Planung und Optimierung von Kurs-
oder Inhaltsabfolgen unter Berücksichtigung fachlicher Abhängigkeiten,
Voraussetzungen und Lernfortschritte \citep{Nabizadeh2020}.

Die einzelnen Ansätze zur Lernpfadpersonalisierung werden durch sehr
unterschiedliche Optimierungs- und Suchverfahren realisiert. Der Survey
von Nabizadeh et al.~zeigt, dass insbesondere evolutionäre Verfahren wie
Genetic Algorithms, schwarmbasierte Optimierungsverfahren wie Ant Colony
Optimization und Particle Swarm Optimization, graphbasierte Verfahren,
Constraint-based Approaches sowie Case-Based-Reasoning-Ansätze
eingesetzt werden \citep{Nabizadeh2020}.

Gemeinsam ist diesen Verfahren, dass sie nicht lediglich einzelne
Inhalte auswählen, sondern mögliche Lernsequenzen als
Optimierungsproblem betrachten. Ziel ist es, unter Berücksichtigung von
Lernzielen, Vorwissen, Voraussetzungen oder zeitlichen Restriktionen
eine möglichst geeignete Reihenfolge von Lernaktivitäten zu bestimmen.
Graphbasierte Verfahren modellieren Lernpfade beispielsweise als
Netzwerke von Lernobjekten und suchen geeignete Wege durch diese
Strukturen. Evolutionäre und schwarmbasierte Verfahren erzeugen dagegen
schrittweise verbesserte Lernpfade, indem verschiedene
Lösungsalternativen miteinander verglichen und optimiert werden
\citep{Nabizadeh2020}.

Darüber hinaus zeigt die Literatur, dass Lernpfadverfahren häufig
mehrere Modellierungsansätze kombinieren. Neben klassischen
Optimierungsverfahren kommen unter anderem
Reinforcement-Learning-Ansätze, Fuzzy-Logic-Methoden sowie
wissensbasierte Verfahren zum Einsatz. Die algorithmische Vielfalt
verdeutlicht, dass Lernpfadpersonalisierung weniger als einzelner
Algorithmus, sondern vielmehr als eigenständige Klasse sequenzieller
Entscheidungs- und Optimierungsprobleme verstanden werden kann
\citep{Nabizadeh2020}.

Insgesamt verdeutlicht die Literatur, dass Lernpfadpersonalisierung
einen eigenständigen Bereich adaptiver Lernverfahren darstellt. Im
Unterschied zu klassischen Recommender-Systemen steht nicht primär die
Auswahl eines einzelnen Inhaltsobjekts im Vordergrund, sondern die
Planung, Generierung und Optimierung ganzer Lernsequenzen
\citep{Nabizadeh2020}. Die Vielfalt der eingesetzten Verfahren -- von
graphbasierten und wissensbasierten Ansätzen über evolutionäre und
schwarmbasierte Optimierungsverfahren bis hin zu
Reinforcement-Learning-Methoden -- zeigt, dass Lernpfadpersonalisierung
als sequenzielles Entscheidungs- und Optimierungsproblem verstanden
werden kann. Für eine informatische Klassifikation adaptiver Verfahren
ergibt sich daraus, dass nicht nur die verwendeten Algorithmen, sondern
auch die Art der zu lösenden Entscheidungsprobleme berücksichtigt werden
sollte.

\subsection{Compurterized Adaptive
Testing}\label{compurterized-adaptive-testing}

Die Grundlage vieler Computerized-Adaptive-Testing-Systeme bilden
psychometrische Messmodelle, insbesondere Verfahren der Item Response
Theory (IRT) \citep{Han2018, Liu2024}. Im Gegensatz zu klassischen Tests
wird dabei nicht davon ausgegangen, dass alle Testpersonen dieselben
Aufgaben bearbeiten müssen. Stattdessen wird die Schwierigkeit einer
Aufgabe in Beziehung zur geschätzten Fähigkeit einer Person gesetzt. Auf
Grundlage der bisherigen Antworten wird fortlaufend eine Schätzung des
Fähigkeitsniveaus berechnet, die anschließend die Auswahl weiterer
Aufgaben steuert.

Han beschreibt die klassische CAT-Architektur über drei zentrale
Komponenten des Item-Selection-Algorithmus: Content Balancing, Item
Selection Criterion und Item Exposure Control \citep{Han2018}. Content
Balancing stellt sicher, dass unterschiedliche Inhaltsbereiche
angemessen berücksichtigt werden. Das Item Selection Criterion bestimmt,
welche Aufgabe für die aktuelle Fähigkeitsschätzung den größten
Informationsgewinn verspricht. Item Exposure Control dient dazu, eine
übermäßige Wiederverwendung einzelner Aufgaben zu verhindern und die
Qualität des Aufgabenpools langfristig zu erhalten.

Die einzelnen Komponenten werden wiederum durch unterschiedliche
psychometrische und algorithmische Verfahren realisiert. Für die
Fähigkeitsschätzung kommen häufig Modelle der Item Response Theory (IRT)
zum Einsatz, insbesondere das Rasch-Modell (1PL) sowie die erweiterten
2PL- und 3PL-Modelle, welche unterschiedliche Eigenschaften von
Testaufgaben berücksichtigen. Die Auswahl der nächsten Aufgabe erfolgt
häufig informationsbasiert, etwa über
Maximum-Fisher-Information-Kriterien, bei denen jene Aufgabe ausgewählt
wird, die für die aktuelle Fähigkeitsschätzung den größten erwarteten
Informationsgewinn liefert. Zur Kontrolle der Aufgabennutzung werden
darüber hinaus Exposure-Control-Verfahren eingesetzt, beispielsweise das
Sympson-Hetter-Verfahren, um eine übermäßige Wiederverwendung einzelner
Testitems zu verhindern \citep{Han2018}.

Der Survey von Liu et al.~erweitert diese klassische Sichtweise um
zusätzliche Komponenten wie Measurement Models, Question Selection, Bank
Construction und Test Control \citep{Liu2024}. Dabei kommen zunehmend
auch Machine-Learning- und Optimierungsverfahren zum Einsatz. Für die
Aufgabenwahl werden beispielsweise informationsbasierte
Selektionsverfahren, Bayes-Methoden oder banditähnliche Strategien
diskutiert. Darüber hinaus werden Verfahren zur automatisierten
Konstruktion und Kalibrierung von Aufgabenpools sowie Modelle zur
Vorhersage von Itemeigenschaften untersucht. Die Literatur verdeutlicht
damit einen Wandel von klassischen psychometrischen CAT-Systemen hin zu
stärker datengetriebenen und algorithmisch erweiterten
Adaptionsarchitekturen \citep{Liu2024}.

Für eine informatische Betrachtung adaptiver Verfahren ist CAT besonders
interessant, weil die einzelnen Schritte der Adaptation explizit
modelliert werden. Die Fähigkeitsschätzung, die Auswahl der nächsten
Aufgabe, die Kontrolle der Testbedingungen sowie die Entscheidung über
die Beendigung des Tests erfolgen innerhalb einer klar definierten
Entscheidungsarchitektur \citep{Han2018, Liu2024}. Im Unterschied zu
vielen anderen adaptiven Lernverfahren steht damit nicht primär die
Modellierung komplexer Lernprozesse im Vordergrund, sondern die
möglichst effiziente Gewinnung von Informationen über den
Fähigkeitsstand einer Person. CAT stellt somit ein Anwendungsfeld dar,
in dem Adaptivität als sequenzielles Entscheidungs- und
Optimierungsproblem besonders deutlich sichtbar wird.

\subsection{Reinforcement Learning}\label{reinforcement-learning}

Reinforcement Learning unterscheidet sich von vielen anderen adaptiven
Verfahren dadurch, dass Entscheidungen nicht auf Grundlage fester Regeln
oder bereits gelernter Modelle getroffen werden, sondern durch
fortlaufende Interaktion mit einer Umgebung optimiert werden. Ein Agent
wählt dabei wiederholt Aktionen aus, beobachtet deren Auswirkungen und
erhält Rückmeldungen in Form von Belohnungen oder Bestrafungen. Ziel ist
es, langfristig eine Strategie zu erlernen, welche den kumulierten
Nutzen über viele Entscheidungsschritte hinweg maximiert
\citep{Riedmann2025}.

Diese Grundidee ist für adaptive Lernsysteme besonders attraktiv, da
viele Bildungsprozesse selbst als sequenzielle Entscheidungsprobleme
verstanden werden können. Die Auswahl eines Lerninhalts beeinflusst
beispielsweise nicht nur den unmittelbaren Lernerfolg, sondern auch
spätere Lernmöglichkeiten und den weiteren Lernverlauf.
Reinforcement-Learning-Verfahren ermöglichen es daher, adaptive
Entscheidungen nicht isoliert, sondern im Kontext langfristiger
Lernziele zu optimieren \citep{Riedmann2025}.

Die Literatur berichtet dabei unterschiedliche Klassen von
Reinforcement-Learning-Verfahren. Zu den klassischen Ansätzen zählt
Q-Learning, bei dem für jede mögliche Aktion schrittweise ein erwarteter
Nutzenwert erlernt wird. Deep-Q-Networks (DQN) erweitern dieses Prinzip
durch neuronale Netze und ermöglichen die Verarbeitung komplexerer
Zustandsräume. Policy-Gradient-Verfahren wie REINFORCE verfolgen dagegen
einen anderen Ansatz: Statt einzelne Aktionswerte zu schätzen, lernen
sie direkt eine Entscheidungsstrategie, welche die Auswahl zukünftiger
Aktionen steuert \citep{Riedmann2025}.

Die verschiedenen Verfahren unterscheiden sich insbesondere hinsichtlich
ihrer Repräsentation von Zuständen, ihrer Lernstrategie sowie ihrer
Skalierbarkeit auf komplexe Lernumgebungen. Während klassische Verfahren
wie Q-Learning häufig leichter interpretierbar sind, ermöglichen
Deep-Reinforcement-Learning-Ansätze die Modellierung deutlich
komplexerer Adaptionsprozesse \citep{Riedmann2025}.

\subsection{Heterogene Verfahren und
Datenquellen}\label{heterogene-verfahren-und-datenquellen}

Über diese Verfahrensfamilien hinaus bleibt festzuhalten, dass adaptive
E-Learning-Systeme in der Praxis nur selten auf einem einzelnen
Modellierungsansatz beruhen. Der systematische Review von Wang et
al.~zeigt, dass insbesondere Lernstile, Wissensstände und Lernverhalten
häufig gemeinsam innerhalb von Learner Models berücksichtigt werden und
zu ihrer Modellierung sehr unterschiedliche Verfahren eingesetzt werden
\citep{Wang2025}.

Die Autoren identifizieren unter anderem Bayesian Networks, Neural
Networks, Fuzzy-Logic-Methoden, Entscheidungsbäume,
k-Nearest-Neighbor-Verfahren, Faktorenanalyse sowie regelbasierte
Ansätze als wiederkehrende Techniken innerhalb adaptiver Lernsysteme
\citep{Wang2025}. Die Bandbreite dieser Verfahren verdeutlicht, dass
adaptive Systeme unterschiedliche Arten von Wissen und Daten verarbeiten
müssen und daher häufig mehrere Modellierungsansätze miteinander
kombinieren.

Besonders interessant ist dabei die Entwicklung der verwendeten
Verfahren über die Zeit. Wang et al.~beschreiben einen Wandel von einer
früher vergleichsweise starken Dominanz von Fuzzy-Logic-Verfahren hin zu
einer deutlich vielfältigeren Methodenlandschaft, in der
probabilistische Modelle, Machine-Learning-Verfahren, neuronale Netze
und regelbasierte Ansätze parallel eingesetzt werden \citep{Wang2025}.
Adaptive Lernsysteme entwickeln sich damit zunehmend zu hybriden
Architekturen, in denen unterschiedliche Verfahren jeweils spezifische
Teilaufgaben übernehmen.

Für eine informatische Klassifikation adaptiver Lernverfahren ergibt
sich daraus eine wichtige Konsequenz: Die einzelnen Verfahrensfamilien
treten in realen Systemen häufig nicht isoliert auf. Ein adaptives
System kann beispielsweise probabilistische Wissensmodelle zur
Lernstandsdiagnostik, Recommender-Verfahren zur Inhaltsauswahl sowie
Optimierungsverfahren zur Lernpfadplanung gleichzeitig verwenden. Die
algorithmische Landschaft adaptiver Lernsysteme ist daher nicht nur
vielfältig, sondern auch stark durch die Kombination unterschiedlicher
Modellierungs- und Entscheidungsverfahren geprägt \citep{Wang2025}.

Neben den verwendeten Modellierungs- und Entscheidungsverfahren
unterscheiden sich adaptive Lernsysteme auch hinsichtlich der
Datenarten, die sie verarbeiten. Die Literatur zeigt, dass
unterschiedliche Verfahrensfamilien auf sehr verschiedenen
Informationsquellen basieren und dadurch jeweils andere Aspekte des
Lernprozesses erfassen \citep{Wang2025, Shen2024, daSilva2023}.

Im Bereich des Knowledge Tracing bilden insbesondere Aufgabenfolgen,
Antwortmuster und Bearbeitungshistorien die Grundlage der
Wissensmodellierung. Erweiterte Modelle berücksichtigen darüber hinaus
zusätzliche Informationen wie Antwortzeiten, Vergessenseffekte,
Lernverhalten oder Tutorinterventionen als sogenannte Side Information
\citep{Shen2024}. Ziel ist es, den Wissenszustand eines Lernenden
möglichst präzise aus einer Folge beobachtbarer Interaktionen
abzuleiten.

Computerized-Adaptive-Testing-Systeme arbeiten ebenfalls mit
Antwortdaten, konzentrieren sich jedoch stärker auf die kontinuierliche
Schätzung von Fähigkeitswerten innerhalb einer Testsituation. Relevant
sind dabei insbesondere Antwortmuster, Itemeigenschaften,
Aufgabenschwierigkeiten sowie Informationen zur Teststeuerung und
Itemauswahl \citep{Han2018, Liu2024}.

Recommender-Systeme greifen dagegen häufig auf deutlich breitere
Datenquellen zurück. Dazu zählen Interessenprofile, Lernhistorien,
Leistungsdaten, Präferenzen, Kompetenzinformationen sowie Eigenschaften
von Lernressourcen. Kollaborative Verfahren nutzen zusätzlich
Interaktionsdaten anderer Lernender, während wissens- und
ontologiebasierte Ansätze stärker auf explizite Domänenmodelle und
strukturierte Wissensrepräsentationen zurückgreifen
\citep{daSilva2023, George2019}.

Auch Machine-Learning-basierte adaptive Lernsysteme verarbeiten häufig
vielfältige Verhaltens- und Kontextdaten. Wang et al.~nennen unter
anderem Lernleistungen, Lernverhalten, demographische Merkmale,
Lernstile sowie Interaktionsdaten als typische Informationsquellen
innerhalb von Learner Models. In cluster- und verhaltensbasierten
Verfahren spielen zusätzlich Klickpfade, Bearbeitungszeiten oder
Nutzungsmuster digitaler Lernumgebungen eine wichtige Rolle
\citep{Wang2025}.

Für eine informatische Klassifikation adaptiver Lernverfahren ergibt
sich daraus, dass Verfahren nicht ausschließlich anhand ihrer
Algorithmen beschrieben werden können. Unterschiedliche
Modellierungsansätze basieren auf unterschiedlichen
Datenrepräsentationen und Informationsquellen. Eine systematische
Einordnung adaptiver Lernverfahren sollte daher neben den verwendeten
Algorithmen auch die zugrunde liegende Datenbasis berücksichtigen, da
diese wesentlich beeinflusst, welche Formen von Adaptivität überhaupt
realisiert werden können.

\section{Zwischenfazit und Anforderungen an eine eigene
Klassifikation}\label{zwischenfazit-und-anforderungen-an-eine-eigene-klassifikation}

Die Analyse des Forschungsstands zeigt, dass adaptive Lernsysteme in der
Literatur aus sehr unterschiedlichen Perspektiven beschrieben und
klassifiziert werden. Während einige Arbeiten Adaptivität primär über
Lernendenmerkmale, pädagogische Zielgrößen oder Formen der
Personalisierung strukturieren, fokussieren andere stärker
Systemarchitekturen, Learner Models, Anwendungskontexte oder konkrete
algorithmische Verfahren. Darüber hinaus werden spezialisierte
Forschungsfelder wie Knowledge Tracing, Educational Recommender Systems,
Computerized Adaptive Testing oder Reinforcement-Learning-basierte
Lernsysteme häufig weitgehend unabhängig voneinander betrachtet, obwohl
sich zahlreiche algorithmische Grundprinzipien zwischen diesen Bereichen
überschneiden.

Die Analyse macht zugleich deutlich, dass bestehende Systematisierungen
häufig unterschiedliche Beschreibungsebenen vermischen. Teilweise stehen
funktionale Rollen adaptiver Verfahren im Vordergrund, etwa
Wissensdiagnostik, Empfehlung oder Sequenzierung. Andere Arbeiten
klassifizieren nach Modellierungsverfahren, Datenquellen, Plattformtypen
oder pädagogischen Zielgrößen. Eine informatisch orientierte
Klassifikation adaptiver Lernverfahren sollte diese unterschiedlichen
Ebenen zunächst analytisch voneinander trennen, bevor sie systematisch
zusammengeführt werden. Erst dadurch lässt sich nachvollziehen, welche
Verfahren ähnliche Probleme lösen, welche Daten sie verarbeiten und
welche Modellierungslogiken ihren Entscheidungen zugrunde liegen.

Aus der Analyse der unterschiedlichen Verfahrensfamilien ergeben sich
mehrere Anforderungen an eine informatisch orientierte Klassifikation
adaptiver Lernverfahren. Die untersuchten Verfahren treten in
unterschiedlichen Forschungskontexten wie Knowledge Tracing, Educational
Recommender Systems, Computerized Adaptive Testing oder
Reinforcement-Learning-basierten Lernumgebungen auf, verwenden jedoch
teilweise ähnliche Modellierungs- und Entscheidungsprinzipien. Eine
Klassifikation sollte daher nicht ausschließlich an Anwendungskontexten
ansetzen, sondern gemeinsame algorithmische Grundideen sichtbar machen.
Gleichzeitig hat die Analyse gezeigt, dass adaptive Verfahren innerhalb
eines Systems unterschiedliche funktionale Rollen übernehmen können,
etwa bei der Wissensdiagnostik, der Empfehlung von Lernressourcen, der
Planung von Lernpfaden oder der Optimierung adaptiver Entscheidungen.
Darüber hinaus unterscheiden sich die Verfahren erheblich hinsichtlich
ihrer Datenquellen, Modellierungsansätze sowie ihrer Transparenz und
Interpretierbarkeit. Eine tragfähige Klassifikation sollte diese
unterschiedlichen Beschreibungsebenen systematisch berücksichtigen.

Zusätzlich macht die Literatur auf mehrere Defizite bestehender
Forschungsansätze aufmerksam. Im Bereich des Learner Modelings wird
wiederholt auf fehlende Standardisierung und uneinheitliche
Beschreibungen von Lernendenmodellen hingewiesen \citep{Boeck2025}.
Educational-Recommender-Systeme werden häufig anhand technischer
Qualitätsmaße wie Genauigkeit oder Präzision bewertet, während ihre
tatsächliche pädagogische Wirksamkeit deutlich seltener untersucht wird
\citep{daSilva2023}. Für Reinforcement-Learning-basierte Lernsysteme
verweist die Literatur zudem auf methodische und empirische
Einschränkungen, da viele Arbeiten weiterhin auf simulierten Szenarien
beruhen und groß angelegte Evaluationen mit realen Lernenden bislang
vergleichsweise selten sind \citep{Riedmann2025}. Diese Befunde
verdeutlichen die weiterhin starke Fragmentierung der Forschung zu
adaptiven Lernverfahren. Eine eigene Klassifikation sollte daher einen
gemeinsamen Ordnungsrahmen schaffen, der Verfahren auch über die Grenzen
einzelner Forschungsfelder hinweg vergleichbar macht.

Für den weiteren Aufbau der Arbeit bedeutet dies, dass Kapitel 2
zunächst die zentralen Perspektiven, Verfahrensfamilien und
Modellierungsansätze adaptiver Lernsysteme systematisch herausarbeitet
und zugleich die Grenzen bestehender Systematisierungen sichtbar macht.
Die Analyse hat gezeigt, dass weder anwendungsbezogene Einteilungen noch
rein algorithmische oder pädagogische Klassifikationen die Vielfalt
adaptiver Lernverfahren vollständig erfassen können. Insbesondere die
Überschneidungen zwischen unterschiedlichen Forschungsfeldern, die
Heterogenität der Datenquellen sowie die Kombination mehrerer Verfahren
innerhalb realer Systeme erschweren eine eindeutige Einordnung
bestehender Ansätze.

Kapitel 3 sollte daher auf dieser Grundlage eine eigenständige
Klassifikation entwickeln, die adaptive Lernverfahren nicht
ausschließlich über Anwendungskontexte oder einzelne Technologien
beschreibt, sondern als informatische Modellierungs-, Diagnose-,
Entscheidungs- und Optimierungsverfahren systematisch einordnet. Der
Mehrwert einer solchen Klassifikation liegt insbesondere darin,
unterschiedliche Verfahrensfamilien über gemeinsame Strukturmerkmale
vergleichbar zu machen und dadurch einen übergreifenden Ordnungsrahmen
für adaptive Lernsysteme bereitzustellen.

\section{Offene Punkte und Grenzen}\label{offene-punkte-und-grenzen}

Die verfügbare Überblicksliteratur zu adaptiven Lernsystemen ist
umfangreich, zugleich jedoch stark fragmentiert. Allgemeine Reviews
fokussieren häufig pädagogische Konzepte, Formen der Personalisierung
oder Systemarchitekturen, während detaillierte Analysen algorithmischer
Verfahren überwiegend in spezialisierten Forschungsfeldern wie Knowledge
Tracing, Educational Recommender Systems, Computerized Adaptive Testing
oder Reinforcement Learning zu finden sind. Eine übergreifende
Betrachtung adaptiver Lernverfahren erfordert daher die Zusammenführung
von Erkenntnissen aus mehreren Teilgebieten, die unterschiedliche
Begrifflichkeiten, Modellierungsansätze und Bewertungskriterien
verwenden.

Darüber hinaus zeigt sich, dass die Evidenzlage häufig stärker auf die
Beschreibung einzelner Verfahren oder Forschungsfelder ausgerichtet ist
als auf deren systematischen Vergleich. Zwar liegen zahlreiche
Untersuchungen zu Lernwirkungen, Vorhersageleistungen oder technischen
Eigenschaften einzelner Ansätze vor, jedoch existieren nur begrenzte
Erkenntnisse darüber, unter welchen Bedingungen bestimmte
Verfahrensfamilien gegenüber anderen vorzuziehen sind. Insbesondere
hybride Systeme, die mehrere Modellierungs- und Entscheidungsverfahren
kombinieren, erschweren eine eindeutige Zuordnung zu bestehenden
Kategorien.

Die im Rahmen dieses Kapitels entwickelte Analyse kann daher keinen
Anspruch auf eine vollständige oder abschließende Bewertung aller
adaptiven Lernverfahren erheben. Sie schafft jedoch eine fundierte
Grundlage, um in den folgenden Kapiteln eine eigenständige informatisch
orientierte Klassifikation zu entwickeln, die die identifizierten
Perspektiven, Verfahrensfamilien und Beschreibungsebenen systematisch
zusammenführt.
